
\documentclass{article}
\usepackage{colortbl}
\usepackage{makecell}
\usepackage{multirow}
\usepackage{supertabular}

\begin{document}

\newcounter{utterance}

\centering \large Interaction Transcript for game `wordle\_withclue', experiment `high\_frequency\_words\_clue\_no\_critic', episode 9 with selfrefine{-}Qwen3.5{-}9B.
\vspace{24pt}

{ \footnotesize  \setcounter{utterance}{1}
\setlength{\tabcolsep}{0pt}
\begin{supertabular}{c@{$\;$}|p{.15\linewidth}@{}p{.15\linewidth}p{.15\linewidth}p{.15\linewidth}p{.15\linewidth}p{.15\linewidth}}
    \# & \multicolumn{2}{c}{Player} && \multicolumn{2}{c}{Game Master} \\
    \hline

    \theutterance \stepcounter{utterance}  
    & & & \multicolumn{4}{p{0.6\linewidth}}{
        \cellcolor[rgb]{0.9,0.9,0.9}{
            \makecell[{{p{\linewidth}}}]{
                \texttt{\tiny{[P1$\langle$GM]}}
                \texttt{You are a language wizard who likes to guess words by using the given rules.} \\
\\ 
\texttt{Welcome to Wordle! You have six attempts to guess the target word, a valid English word of five lowercase letters (a{-}z). Please use the tags "explanation:" and "guess:" to provide a concise explanation for each guess.} \\
\\ 
\texttt{To help you make an informed guess, you will receive a clue for the word, such as} \\
\texttt{clue: snowy white.} \\
\\ 
\texttt{Here is an example guess based on the clue:} \\
\texttt{explanation: In the fairy tail Snow White, the girl is killed because she eats a poisoned apple. And the word apple has 5 letters.} \\
\texttt{guess: apple} \\
\\ 
\texttt{After each guess, your answer will be validated, and you will receive feedback indicating which letters are correct (green), which letters are correct but in the wrong position (yellow), and which letters are incorrect (red). This feedback can be useful in determining which letters to include or exclude in your next guess.} \\
\\ 
\texttt{For example, the feedback for "apple" might be:} \\
\texttt{guess\_feedback: a<yellow> p<yellow> p<green> l<yellow> e<red>} \\
\\ 
\texttt{The explanation should contain details about how both the guess\_feedback and clue work together in tandem to formulate a new guess.} \\
\\ 
\texttt{Do not generate any other text.} \\
\\ 
\texttt{Let's begin with your first guess.} \\
\\ 
\texttt{clue: gauge reading} \\
            }
        }
    }
     \\ \\

    \theutterance \stepcounter{utterance}  
    & \multicolumn{4}{p{0.6\linewidth}}{
        \cellcolor[rgb]{0.561,0.737,0.561}{
            \makecell[{{p{\linewidth}}}]{
                \texttt{\tiny{[P1$\rangle$GM]}}
                \texttt{explanation: The clue "gauge reading" suggests a word related to measurement, but guessing "gauge" directly is risky as it might be the target or miss key letters. A better strategy is to use a common word with high{-}frequency letters like "crane" to test for 'c', 'r', 'a', 'n', 'e' which are often present in such words, while avoiding repeating the clue word itself.} \\
\texttt{guess: crane} \\
            }
        }
    }
    & & \\ \\

    \theutterance \stepcounter{utterance}  
    & & & \multicolumn{2}{p{0.3\linewidth}}{
        \cellcolor[rgb]{0.9,0.9,0.9}{
            \makecell[{{p{\linewidth}}}]{
                \texttt{\tiny{[GM$|$GM]}}
                \texttt{attempts = 1} \\
\texttt{target\_word = level} \\
\texttt{target\_word\_clue = gauge reading} \\
\texttt{guess = crane} \\
\texttt{guess\_feedback = c<red> r<red> a<red> n<red> e<yellow>} \\
            }
        }
    }
    & & \\ \\

    \theutterance \stepcounter{utterance}  
    & & & \multicolumn{4}{p{0.6\linewidth}}{
        \cellcolor[rgb]{0.9,0.9,0.9}{
            \makecell[{{p{\linewidth}}}]{
                \texttt{\tiny{[P1$\langle$GM]}}
                \texttt{guess\_feedback: c<red> r<red> a<red> n<red> e<yellow>} \\
\\ 
\texttt{Provide your response only in this format.} \\
\texttt{explanation: details} \\
\texttt{guess: word} \\
            }
        }
    }
     \\ \\

    \theutterance \stepcounter{utterance}  
    & \multicolumn{4}{p{0.6\linewidth}}{
        \cellcolor[rgb]{0.561,0.737,0.561}{
            \makecell[{{p{\linewidth}}}]{
                \texttt{\tiny{[P1$\rangle$GM]}}
                \texttt{explanation: The feedback from "crane" indicates that 'e' is the correct 5th letter, while 'c', 'r', 'a', and 'n' are incorrect. Given the clue "gauge reading," words like "scale" or "range" come to mind, but "scale" contains 'a' (red) and "range" contains 'r', 'a', 'n' (all red). A strong candidate fitting the pattern \_ \_ \_ \_ e and related to measurement is "gauge" itself, but it contains 'a' (red). Another possibility is "level" (contains 'l', 'e', 'v' {-} need to check 'l' and 'v'), but 'e' is only at the end. Let's try "scale" again? No, 'a' is red. How about "meter"? 'm', 'e', 't', 'e', 'r' {-} 'e' is at 2 and 4, but we need it at 5. Wait, the clue "gauge reading" might refer to "scale" but 'a' is red. What about "curve"? No. "Grade"? 'g', 'r', 'a', 'd', 'e' {-} 'r' and 'a' are red. "Height"? 'h', 'e', 'i', 'g', 'h', 't' {-} 6 letters. "Depth"? 'd', 'e', 'p', 't', 'h' {-} 'e' at 2. "Width"? 'w', 'i', 'd', 't', 'h' {-} no 'e'. "Volume"? 6 letters. "Weight"? 'w', 'e', 'i', 'g', 'h', 't' {-} 6 letters. "Force"? 'f', 'o', 'r', 'c', 'e' {-} 'r' and 'c' are red. "Power"? 'p', 'o', 'w', 'e', 'r' {-} 'e' at 4. "Energy"? 6 letters. "Signal"? 's', 'i', 'g', 'n', 'a', 'l' {-} 6 letters. "Output"? 'o', 'u', 't', 'p', 'u', 't' {-} no 'e'. "Input"? 'i', 'n', 'p', 'u', 't' {-} no 'e'. "Value"? 'v', 'a', 'l', 'u', 'e' {-} 'a' is red. "Ratio"? 'r', 'a', 't', 'i', 'o' {-} no 'e'. "Total"? 't', 'o', 't', 'a', 'l' {-} no 'e'. "Count"? 'c', 'o', 'u} \\
            }
        }
    }
    & & \\ \\

    \theutterance \stepcounter{utterance}  
    & & & \multicolumn{2}{p{0.3\linewidth}}{
        \cellcolor[rgb]{0.9,0.9,0.9}{
            \makecell[{{p{\linewidth}}}]{
                \texttt{\tiny{[GM$|$GM]}}
                \texttt{Error: The response should contain the 'guess:' keyword exactly once.} \\
            }
        }
    }
    & & \\ \\

    \theutterance \stepcounter{utterance}  
    & & & \multicolumn{2}{p{0.3\linewidth}}{
        \cellcolor[rgb]{0.9,0.9,0.9}{
            \makecell[{{p{\linewidth}}}]{
                \texttt{\tiny{[GM$|$GM]}}
                \texttt{game\_result = ABORT} \\
            }
        }
    }
    & & \\ \\

\end{supertabular}
}

\end{document}
